\begin{mistake}{2026-04-09}{02}

\begin{problemcontent}
设常数 \( a > 0 \)，则 \( \lim_{x \to +\infty} \frac{a^x}{1+a^x} = \) \underline{\hspace{2em}}。
\end{problemcontent}

\begin{solutioncontent}
指数函数 \(a^x\) 当 \(x \to +\infty\) 时的极限行为取决于底数 \(a\) 与 \(1\) 的大小关系，需分类讨论：

\begin{enumerate}
 \item 当 \(0 < a < 1\) 时，\(\lim_{x \to +\infty} a^x = 0\)，原极限为 \(\frac{0}{1+0} = 0\)。
 \item 当 \(a = 1\) 时，\(a^x \equiv 1\)，原极限为 \(\frac{1}{1+1} = \frac{1}{2}\)。
 \item 当 \(a > 1\) 时，\(\lim_{x \to +\infty} a^x = +\infty\)，分子分母同除以 \(a^x\)：
 \[\lim_{x \to +\infty} \frac{1}{a^{-x} + 1} = \frac{1}{0 + 1} = 1\]
\end{enumerate}
\end{solutioncontent}

\mistakeknowledge{
\begin{itemize}
 \item \textbf{核心定理}：指数函数 \(y=a^x\) 的极限性质及渐近线。
 \item \textbf{易错点}：忽略底数 \(a\) 的范围，未分类讨论而直接得出 \(1\) 的错误结论。
 \item \textbf{关联知识}：当自变量趋于无穷时，必须区分 \(x \to +\infty\) 与 \(x \to -\infty\) 对 \(a^x\) 极限截然不同的影响。
\end{itemize}}

\end{mistake}
